%% ==========================================================================
%% SECTION 1: INTRODUCTION (Paxos framing)
%% ==========================================================================
\section{Introduction}\label{sec:intro}

Systems that incorporate language model components face a verification problem
with a precise structural analogy in distributed systems. In Paxos
consensus~\cite{paxos}, three roles divide the work: \emph{proposers} generate
candidate values, \emph{acceptors} filter proposals into decisions, and
\emph{learners} observe the decided values. A learner that sees only the decided
sequence cannot reconstruct what was proposed and rejected---the decision
process is opaque by design. We show that language model generation exhibits the
same structure, and that the opacity creates the same class of verification
problem.

In a transformer, the internal computation---feature activations, competing
attention heads, probability mass spread across the vocabulary---plays the role
of proposer: many candidate continuations are represented in parallel. Post-training
alignment (RLHF, HHA) plays the role of acceptor: it biases selection toward
tokens that satisfy helpfulness, harmlessness, and coherence objectives.
The output token stream is the decided sequence. External
observers---users, supervisors, downstream system components---are learners.
They see only the linearized output. They cannot see what was proposed and
rejected, how contested each decision was, or whether the acceptor overrode a
high-entropy proposal distribution to emit a confident-sounding token.
The token stream is a projection of a parallel computation, and the
projection discards the epistemic state that verification requires.

\fakepara{The problem.}
This opacity is not incidental. It is architectural, and it has direct
consequences for verification cost. A supervisor receiving only text---only the
decided values---must spend its verification budget reconstructing what the
model's own computation already knew. The budget is consumed by
reconstruction rather than concentrated on the cases where verification is most
needed. We prefer the term \emph{fabrication} to emphasize that fluent
confabulation is the system's default behavior under coherence optimization, not
a malfunction to be patched. Xu et al.~\cite{xu2024hallucination} prove a
related result: hallucination is inevitable for LLMs over the class of
computable functions. Our result is orthogonal. They prove a computational limit
on what the model can produce; we prove an observational limit on what a
supervisor can verify. The model may know it is fabricating. The supervisor
cannot tell.

\fakepara{Why self-report fails.}
The Paxos analogy predicts a specific failure mode, and the empirical data
confirms it. In consensus, a decided value carries no trace of how contested the
vote was. In language models, the acceptor (alignment training) optimizes for
helpfulness and fluency, which means the decided token often expresses
confidence regardless of how uncertain the proposal distribution was. Across
four model architectures (OLMo-3 7B, Llama-3.1 8B, Qwen3 4B, Mistral 7B),
every model reports \emph{higher} confidence on fabrications than on known
facts. Self-report AUC ranges 0.28--0.36---below random, because the
relationship is systematically inverted. Self-report inversion is not a training
failure. It is the acceptor working as designed: asked ``how confident are
you?'', the helpful-harmless-honest filter proposes a confident answer, because
confident answers score well on helpfulness. The proposer may have been
uncertain. The learner sees only what the acceptor decided.

\fakepara{The structural challenge.}
Stacking learners does not help. A second language model judging the first
model's output is another learner observing the same decided sequence. It sees
no finer information than the original observer. Text-only supervision is
closed under composition: judges of judges remain in the same observation class.
This is the distributed systems insight: adding more learners to a Paxos
deployment does not give any learner access to the proposal history. The
information was discarded at the interface boundary, and no downstream
processing can recover it.

There is a deeper structural problem. In standard Paxos, acceptor behavior is
fixed by the protocol: the acceptance criterion does not change between rounds.
In a language model, the acceptor redefines its own weights each round. Each
output token changes the context, which changes attention patterns, which
changes the effective acceptance criterion for the next token. The model is
simultaneously a consensus protocol and a dynamically reconfiguring one. This is
what makes language models both capable and unreliable: the adaptive acceptor
enables flexible, context-sensitive generation, but it also means that the
proposer--acceptor relationship is non-stationary. A verification strategy that
works for one region of the generation may fail for the next.

\fakepara{Key insight.}
The generation process produces signals that the text interface
discards---\emph{telemetry} (measurements of the computation that produced the
output) as distinct from \emph{testimony} (what the model says). The entropy of
the proposal distribution at each step records how contested the decision was.
Attention summaries record which context positions influenced the acceptor. These
are proposer-level signals: they describe the state of the computation
\emph{before} acceptance collapsed it into a single token. Under current
training objectives, the model cannot independently control this telemetry---it
is a byproduct of the computation, not a product of the optimization objective.
The testimony may lie; the telemetry is harder to fake.

\fakepara{Our approach.}
We construct a \emph{tensor interface} that exports proposer-level telemetry to
learners, escaping the observation class in which the impossibility holds
(\autoref{sec:escape}). The interface requires no architectural changes to
current transformers. It exports entropy traces and attention summaries alongside
the token stream, enabling external supervisors to triangulate: tensor signals
read how wide the proposal distribution was; text signals read what survived
acceptance. Neither alone can distinguish ``wide proposals, stable acceptance''
(genuine uncertainty resolved by context) from ``narrow proposals, unstable
acceptance'' (confident computation with a drifting acceptor)---but composed
together, they can.

\fakepara{Results.}
Under bounded verification budgets, tensor-guided judges at 10\% budget
outperform text-guided judges at 30\% budget (82.1\% vs.\ 78.5\% end-to-end
accuracy, averaged across four architectures): the tensor enables a $3\times$
reduction in verification cost. The text channel is stronger than
expected---response length alone achieves AUC 0.88--0.96---but it is trivially
gameable: a model trained to fabricate briefly defeats length-based detection.
After controlling for length, the tensor residual adds AUC 0.60--0.70: a
computational signal that cannot be independently gamed without adversarial
optimization targeting the signal itself. The discrimination is architectural,
not model-specific: cross-model agreement reaches Spearman $\rho = 0.762$,
meaning the proposer-level signal generalizes across model families.

\fakepara{Contributions.}
This paper establishes three results that together characterize the cost surface
of epistemic verification.

\begin{itemize}
\item \textbf{An impossibility result as budget constraint.}
  \autoref{sec:formal} proves that epistemic honesty is unverifiable for
  predictive token generation systems operating under text-only observation with
  bounded supervision. The impossibility is structural: it holds regardless of
  model scale, training procedure, or prompt engineering, and cannot be escaped
  by stacking supervisors---just as adding Paxos learners does not recover
  discarded proposal history. It establishes the floor below which text-channel
  verification investment cannot improve assurance.

\item \textbf{A tensor interface as the next investment tier.}
  \autoref{sec:escape} demonstrates a minimal interface extension that exports
  proposer-level telemetry alongside text with no architectural changes to
  current transformers. The tensor interface increases the observation budget
  available to external verifiers by exposing computational signals the
  text-only interface discards.

\item \textbf{An empirical cost surface.}
  \autoref{sec:eval} decomposes verification effectiveness into text-channel and
  tensor-channel components, characterizes the cost-robustness tradeoff, and
  identifies the format-dependent variation that determines optimal judge
  strategy per domain. The contribution is the map, not a single detector: systems
  engineers building with language model components need to know what each level
  of verification investment buys. This paper provides the cost surface for that
  decision.
\end{itemize}
